\scanpage{76}
Nous nous intéressons au cas $n=2$, et au cas où $G_0$ est fini et \deleted{\unsure{admet une représentation}}\inserted{a sa représentation en $\operatorname{Aut}(X)$} \uline{fidèle} \deleted{\unsure{comme dans} $\operatorname{Aut}(X)$} (peut-être pas dans $\operatorname{Aut}(C)$. -- il peut y avoir un noyau non trivial pour $G_0\unsure{\text{---}|\longrightarrow}\operatorname{Aut}(C)$]

\[
 \left\{\begin{aligned}
 &\text{Invariants}\qquad\beta_0=\alpha_0+2,\quad\beta_1=\alpha_1+2\\
 &\rho=8-2(\beta_0+\beta_1)=-2(\alpha_0+\alpha_1)\qquad\widetilde\rho=-(\alpha_0+\alpha_1)\\
 &\sigma=4-\beta_1=2-\alpha_1,\quad\sigma'=4-\beta_0=2-\alpha_0
 \end{aligned}\right.
\]

$Q_C=Y$ de dim relative $1$, lisse sss $\beta_0\beta_1\widetilde\rho$ inv, i.e. sss en tout point on a \deleted{\unsure{$\alpha_0$}}\inserted{$\alpha_0\ne-2,\alpha_1\ne-2,$} $\alpha_0+\alpha_1\ne0$

$Q_X$ de dim relative $2$, lisse sss $\beta_1\beta_1\sigma$ inv, i.e. sss en tout point on a

\[
 \alpha_0\ne-2,\quad\alpha_1\ne\pm2,
\]

\[
\begin{aligned}
 \rho_f&=\tau_0\tau_1\quad\text{son ordre dépend de la façon connue de }\beta_0=\alpha_0+2\\
 \rho_s&=\tau_1\tau_2\hspace{46mm}\beta_1=\alpha_1+2,
\end{aligned}
\]

\deleted{\unsure{D'autre|D'où}} \unsure{On} sait que

% slanted
\margin{vérifier si \unsure{$Q$|$\varnothing$} non lisse}

\[
 \underline{\operatorname{Aut}}(C,Q_C)\overset{\sim}{\longrightarrow}
 \underline{\operatorname{Aut}}(Q_C)=\operatorname{Aut}(Y)
\]

Donc la cat. des 2-polyèdres réguliers \deleted{\unsure{et sur $S$}}\inserted{\unsure{paramétrés}} épinglés sur $S$ \inserted{rel. à $G_0$} est équivalente à celle des triples $(Y,\varphi_{G_0}^{Y})$ d'une \unsure{cônique|conique} $Y$ sur $S$ et d'un hom.

% slanted
\margin{$G_0$ fini ou non, mais il faut supposer $(\beta_0,2)\mathcal O_S=\mathcal O_S$ i.e. $\beta_0\ne0$ aux pts de car. $2$}

\[
 \varphi_{G_0}^{Y}:G_0\longrightarrow\operatorname{Aut}(Y),
\]

satisfaisant les conditions

a) \deleted{\unsure{$\tau_0,\tau_1,\tau_2$}} $\rho_s^2\ne1,\ \rho_f^2\ne1$ en \unsure{toute|toutes} fibre ($\Rightarrow\tau_0,\tau_1,\tau_2\ne\mathrm{id}$ sur \unsure{tte} fibre, donc $\tau_0,\tau_1,\tau_2$ des \uline{réflexions})

% slanted
\margin{$2$ \deleted{\unsure{pour que}} a) (avec b) \deleted{soit} \uline{suffisant} pour que \deleted{ça provienne d'une} structure de polyèdre régulier, il faut se limiter au cas où $\deleted{\unsure{\rho=}}\beta_0\beta_1$ inv (les $\tau_i^C$ sont des réflexions \deleted{\unsure{affines}} \deleted{\unsure{droites}} orthogonales)}

b) Les sections $h_0,h_1,h_2$ de $C$ définies par $\tau_0,\tau_1,\tau_2$ sont non alignées en tout point. \deleted{Si}

Si $\beta_0$ inv et $(\beta_1,4-\beta_0)\mathcal O_S=\mathcal O_S$ [i.e. \deleted{\unsure{$\beta_1$, en un pt, on a}} \inserted{en aucun pt on n'a} à la fois \deleted{\unsure{$\beta_1=0,\beta$}} $\alpha_1=-2,\alpha_0=2$], alors b) équivaut à

b') En tout pt, la repr. $\varphi_{G_0}^C$ \deleted{\unsure{est}} a pas de droite fixe.

Si $Y$ lisse, i.e. \deleted{\unsure{$\beta_0,\beta_1,\widetilde\rho$, de}}\inserted{en aucun} pt on n'a \deleted{\unsure{simultanément}}\inserted{l'une des égalités} $\alpha_0=-2,\alpha_1=-2,\alpha_0+\alpha_1=0$, alors la condition b') équivaut à :

b'') \inserted{Sur aucune fibre} Il n'existe \deleted{pas} de $2$-diviseur stable pour $\varphi_{G_0}^{Y}$.
% archive sheet number: 75
